\begin{table}[htbp]\centering\small
\caption{Salud: los vehículos presupuestales, línea P (mdp)}\label{cua:salud}
\begin{tabular}{lrrr}
\toprule
Vehículo & 2026 & 2027 & Var. real (\%) \\
\midrule
Ramo 12 Salud & 66.825,8 & 72.305,5 & +4,0 \\
Ramo 56 IMSS-Bienestar & 172.492,4 & 193.987,2 & +8,1 \\
IMSS (entidad de control directo) & 1.590.308,7 & 1.803.009,1 & +9,0 \\
ISSSTE (entidad de control directo) & 539.020,8 & 571.592,3 & +2,0 \\
FASSA (Ramo 33) & 84.635,9 & 89.714,0 & +1,9 \\
\bottomrule
\end{tabular}
\\[2pt]\begin{minipage}{0.92\textwidth}\scriptsize \textbf{Estos cinco renglones NO se suman}: el IMSS y el ISSSTE incluyen pensiones y prestaciones que no son atención médica, y el Ramo 19 les transfiere recursos que volverían a contarse. Fuentes: Anexo 1 y Anexo 22 de los dos proyectos de decreto.\end{minipage}
\end{table}
